% Supplementary F — Journal §VI (Results)
% Extended per-axis results: full tables for F-1..F-12 across 6 scenes
% + 5 HIDSAG subsets.

\documentclass[11pt,a4paper]{article}
% Shared preamble for all supplementary chapter documents.
% Used by each supplementary/conference/suppl_*.tex and
% supplementary/journal/suppl_*.tex via % Shared preamble for all supplementary chapter documents.
% Used by each supplementary/conference/suppl_*.tex and
% supplementary/journal/suppl_*.tex via % Shared preamble for all supplementary chapter documents.
% Used by each supplementary/conference/suppl_*.tex and
% supplementary/journal/suppl_*.tex via \input{../preamble.tex}.

\usepackage[utf8]{inputenc}
\usepackage[T1]{fontenc}
\usepackage[a4paper,margin=2.2cm]{geometry}
\usepackage{amsmath,amssymb,amsfonts,amsthm}
\usepackage{mathtools}
\usepackage{graphicx}
\usepackage{booktabs}
\usepackage{array}
\usepackage{multirow}
\usepackage{longtable}
\usepackage{algorithm}
\usepackage{algpseudocode}
\usepackage{xcolor}
\usepackage{url}
\usepackage{cite}
\usepackage[hidelinks]{hyperref}

\graphicspath{{../../figures/}}

\renewcommand{\thesection}{\Alph{section}.\arabic{section}}
\renewcommand{\thesubsection}{\thesection.\arabic{subsection}}

\theoremstyle{plain}
\newtheorem{theorem}{Theorem}[section]
\newtheorem{lemma}[theorem]{Lemma}
\newtheorem{proposition}[theorem]{Proposition}
\theoremstyle{definition}
\newtheorem{definition}[theorem]{Definition}
\theoremstyle{remark}
\newtheorem{remark}[theorem]{Remark}

\setcounter{secnumdepth}{3}
\setcounter{tocdepth}{2}

\newcommand{\R}{\mathbb{R}}
\newcommand{\E}{\mathbb{E}}
\newcommand{\KL}[2]{\mathrm{KL}\!\left(#1 \,\middle\|\, #2\right)}
\newcommand{\ari}{\mathrm{ARI}}
\newcommand{\nmi}{\mathrm{NMI}}
\newcommand{\softmax}{\mathrm{softmax}}
.

\usepackage[utf8]{inputenc}
\usepackage[T1]{fontenc}
\usepackage[a4paper,margin=2.2cm]{geometry}
\usepackage{amsmath,amssymb,amsfonts,amsthm}
\usepackage{mathtools}
\usepackage{graphicx}
\usepackage{booktabs}
\usepackage{array}
\usepackage{multirow}
\usepackage{longtable}
\usepackage{algorithm}
\usepackage{algpseudocode}
\usepackage{xcolor}
\usepackage{url}
\usepackage{cite}
\usepackage[hidelinks]{hyperref}

\graphicspath{{../../figures/}}

\renewcommand{\thesection}{\Alph{section}.\arabic{section}}
\renewcommand{\thesubsection}{\thesection.\arabic{subsection}}

\theoremstyle{plain}
\newtheorem{theorem}{Theorem}[section]
\newtheorem{lemma}[theorem]{Lemma}
\newtheorem{proposition}[theorem]{Proposition}
\theoremstyle{definition}
\newtheorem{definition}[theorem]{Definition}
\theoremstyle{remark}
\newtheorem{remark}[theorem]{Remark}

\setcounter{secnumdepth}{3}
\setcounter{tocdepth}{2}

\newcommand{\R}{\mathbb{R}}
\newcommand{\E}{\mathbb{E}}
\newcommand{\KL}[2]{\mathrm{KL}\!\left(#1 \,\middle\|\, #2\right)}
\newcommand{\ari}{\mathrm{ARI}}
\newcommand{\nmi}{\mathrm{NMI}}
\newcommand{\softmax}{\mathrm{softmax}}
.

\usepackage[utf8]{inputenc}
\usepackage[T1]{fontenc}
\usepackage[a4paper,margin=2.2cm]{geometry}
\usepackage{amsmath,amssymb,amsfonts,amsthm}
\usepackage{mathtools}
\usepackage{graphicx}
\usepackage{booktabs}
\usepackage{array}
\usepackage{multirow}
\usepackage{longtable}
\usepackage{algorithm}
\usepackage{algpseudocode}
\usepackage{xcolor}
\usepackage{url}
\usepackage{cite}
\usepackage[hidelinks]{hyperref}

\graphicspath{{../../figures/}}

\renewcommand{\thesection}{\Alph{section}.\arabic{section}}
\renewcommand{\thesubsection}{\thesection.\arabic{subsection}}

\theoremstyle{plain}
\newtheorem{theorem}{Theorem}[section]
\newtheorem{lemma}[theorem]{Lemma}
\newtheorem{proposition}[theorem]{Proposition}
\theoremstyle{definition}
\newtheorem{definition}[theorem]{Definition}
\theoremstyle{remark}
\newtheorem{remark}[theorem]{Remark}

\setcounter{secnumdepth}{3}
\setcounter{tocdepth}{2}

\newcommand{\R}{\mathbb{R}}
\newcommand{\E}{\mathbb{E}}
\newcommand{\KL}[2]{\mathrm{KL}\!\left(#1 \,\middle\|\, #2\right)}
\newcommand{\ari}{\mathrm{ARI}}
\newcommand{\nmi}{\mathrm{NMI}}
\newcommand{\softmax}{\mathrm{softmax}}


\title{Journal Supplementary F --- Extended per-axis results\\
{\large Companion to: \emph{Beyond Accuracy} (Section VI)}}
\author{Felipe Santib\'a\~nez-Leal}
\date{2026}

\begin{document}
\maketitle

\section{Scope}
This supplement provides the full per-axis result tables for the 6
labelled scenes and 5 HIDSAG subsets that did not fit in journal
Section~VI. All numbers are read directly from the public artefacts
in \texttt{CAOS\_LDA\_HSI/data/derived/}.

\section{Axis F-1: hierarchical Bayesian classification --- pairwise $P[A>B]$}

\subsection{Labelled-scene block}

\begin{table}[h]
\caption{Pairwise $P[A > B]$ from the labelled-scene posterior (rows
= method A, columns = method B). Estimated from 2000 NUTS draws.}
\label{tab:b1-pair-labelled}
\centering
\small
\begin{tabular}{l c c c c c}
\toprule
A \textbackslash{} B & pca\_K & raw & theta & hard & soft \\
\midrule
pca\_K\_logistic    & ---   & 0.00 & 1.00 & 0.00  & 0.00  \\
raw\_logistic       & 1.00 & ---   & 1.00 & 0.54  & 0.36  \\
theta\_logistic     & 0.00 & 0.00 & ---   & 0.00  & 0.00  \\
topic\_routed\_hard & 1.00 & 0.46 & 1.00 & ---    & 0.32  \\
topic\_routed\_soft & 1.00 & 0.64 & 1.00 & 0.68   & ---    \\
\bottomrule
\end{tabular}
\end{table}

The dominant observation is in the bottom-right $2 \times 2$ block:
$P[\text{soft} > \text{raw}] = 0.64$, $P[\text{hard} > \text{raw}] =
0.46$, $P[\text{soft} > \text{hard}] = 0.68$. Topic-routed-soft is
the leading method by a small margin; topic-routed-hard and
raw-logistic are statistically indistinguishable. theta\_logistic
alone is dominated by every other method, confirming that
$\theta$-as-feature without the gating wrapper discards too much
information.

\subsection{HIDSAG block}

\begin{table}[h]
\caption{Pairwise $P[A > B]$ from the HIDSAG posterior. Methods:
cube = cube-topic logistic, pca = PCA logistic, raw = raw logistic,
rgn = region-topic logistic, $\theta$ = topic logistic.}
\label{tab:b1-pair-hidsag}
\centering
\small
\begin{tabular}{l c c c c c}
\toprule
A \textbackslash{} B & cube & pca & raw & rgn & $\theta$ \\
\midrule
cube\_topic   & ---   & 0.00 & 0.00 & 0.23 & 0.06 \\
pca\_logistic & 1.00 & ---   & 0.28 & 1.00 & 1.00 \\
raw\_logistic & 1.00 & 0.72 & ---   & 1.00 & 1.00 \\
region\_topic & 0.77 & 0.00 & 0.00 & ---   & 0.19 \\
theta\_logistic & 0.94 & 0.00 & 0.00 & 0.81 & ---   \\
\bottomrule
\end{tabular}
\end{table}

The HIDSAG inversion is the row-pattern of raw\_logistic: 1.00, 0.72,
1.00, 1.00. Raw-spectrum logistic dominates every topic-based method
with $P = 1.00$ against $\theta$-logistic and region-topic. The
ordering is exactly the labelled-scene ordering inverted.

\section{Axis F-2: per-scene coherence}

\begin{table}[h]
\caption{Canonical-fit coherence per scene at $N = 15$ top words.}
\label{tab:b2}
\centering
\small
\begin{tabular}{l c c c}
\toprule
Scene            & $c_v$ & $c_{\text{NPMI}}$ & U-Mass \\
\midrule
Indian Pines     & 0.450 & $-0.646$ & 0.0 \\
Salinas          & ---    & ---       & --- \\
Salinas-A        & 0.409 & ---       & --- \\
Pavia U          & 0.637 & ---       & --- \\
KSC              & 0.556 & ---       & --- \\
Botswana         & 0.473 & ---       & --- \\
\bottomrule
\end{tabular}
\end{table}

\section{Axis F-3: seed stability per (scene, method)}

\begin{table}[h]
\caption{ARI mean $\pm$ std over 5 seeds. $c_v$ mean reported alongside.}
\label{tab:b3-full}
\centering
\small
\begin{tabular}{l c c c c}
\toprule
Scene & Method & ARI mean & ARI std & $c_v$ mean \\
\midrule
\multirow{2}{*}{Indian Pines} & ProdLDA & 0.165 & 0.003 & 0.630 \\
                              & ETM     & 0.257 & 0.003 & 0.357 \\
\addlinespace
\multirow{2}{*}{Salinas}      & ProdLDA & 0.306 & 0.008 & 0.616 \\
                              & ETM     & 0.456 & 0.030 & 0.354 \\
\addlinespace
\multirow{2}{*}{Salinas-A}    & ProdLDA & 0.552 & 0.016 & 0.465 \\
                              & ETM     & 0.584 & 0.018 & 0.351 \\
\addlinespace
\multirow{2}{*}{Pavia U}      & ProdLDA & 0.413 & 0.016 & 0.786 \\
                              & ETM     & 0.458 & 0.025 & 0.694 \\
\addlinespace
\multirow{2}{*}{KSC}          & ProdLDA & 0.195 & 0.018 & 0.654 \\
                              & ETM     & 0.209 & 0.009 & 0.603 \\
\addlinespace
\multirow{2}{*}{Botswana}     & ProdLDA & 0.169 & 0.005 & 0.780 \\
                              & ETM     & 0.267 & 0.007 & 0.412 \\
\bottomrule
\end{tabular}
\end{table}

Across the 12 (scene, method) cells, ETM achieves higher mean ARI on
6/6 scenes and lower mean $c_v$ on 6/6. The within-seed standard
deviation is below 0.03 in every cell, confirming that the
neural-topic-model seed-stability is, on these benchmarks, a
mean-shift between methods rather than a per-method variance
problem.

\section{Axis F-4: capacity sweep --- recommended $K$ per scene}

\begin{table}[h]
\caption{Per-scene recommended $K$ from the LDA sweep. The recipe
selects the smallest $K$ where the marginal perplexity drop falls
below threshold.}
\label{tab:b4-recK}
\centering
\small
\begin{tabular}{l c c c}
\toprule
Scene & Recommended $K$ & $K$ grid & Recipe \\
\midrule
Indian Pines & 4 & \{4,6,8,10,12,16\} & matched-cosine plateau \\
Salinas      & 4 & same & matched-cosine plateau \\
Salinas-A    & 4 & same & matched-cosine plateau \\
Pavia U      & 4 & same & matched-cosine plateau \\
KSC          & 4 & same & matched-cosine plateau \\
Botswana     & 4 & same & matched-cosine plateau \\
\bottomrule
\end{tabular}
\end{table}

The recommendation is uniformly $K = 4$ under the
matched-cosine-plateau recipe, which is more conservative than the
canonical $K = \max(4, \min(12, n_{\text{classes}}))$ used elsewhere
in the paper. The difference is by design: the canonical $K$ is
chosen for downstream classification surface coverage, the
recommended $K$ is chosen for basis-stability.

\section{Axis F-5: full paired ARI per (scene, mask)}

Restated from conference supplement D for completeness: paired ARI
values are in conference table I of the main paper. The full table
with swap rates, bands kept, perplexity, and ARI-vs-label is in
conference supplement D.

\section{Axis F-6: cross-method ARI per scene}

\begin{table}[h]
\caption{Cross-method ARI on Indian Pines (excerpt from journal Table
III). Full per-scene tables in companion artefact
\texttt{cross\_method\_agreement/<scene>.json}.}
\label{tab:b6-ip}
\centering
\small
\begin{tabular}{l c c c c c c}
\toprule
              & lbl & top  & felz  & SLIC500 & SLIC2k & p15 \\
\midrule
label         & ---  & 0.20 & 0.57 & 0.23 & 0.05 & 0.28 \\
topic-dom     & 0.20 & ---  & 0.12 & 0.04 & 0.01 & 0.05 \\
felzenszwalb  & 0.57 & 0.12 & ---  & 0.31 & 0.08 & 0.37 \\
SLIC-500      & 0.23 & 0.04 & 0.31 & ---  & 0.24 & 0.36 \\
SLIC-2000     & 0.05 & 0.01 & 0.08 & 0.24 & ---  & 0.15 \\
patch-15      & 0.28 & 0.05 & 0.37 & 0.36 & 0.15 & ---  \\
\bottomrule
\end{tabular}
\end{table}

\section{Axis F-7: topic--label coupling per scene}

\begin{table}[h]
\caption{Per-scene mean and max KL between $P(L \mid t)$ and the
marginal $P(L)$ on the canonical fit.}
\label{tab:b7}
\centering
\small
\begin{tabular}{l c c}
\toprule
Scene & KL mean (nats) & KL max (nats) \\
\midrule
Indian Pines & 4.53  & 27.33 \\
Salinas      & 7.11  & 27.63 \\
Salinas-A    & 4.79  & 26.23 \\
Pavia U      & 4.09  & 11.53 \\
KSC          & 7.85  & 27.63 \\
Botswana     & 10.48 & 27.63 \\
\bottomrule
\end{tabular}
\end{table}

\section{Axis F-9: HIDSAG preprocessing stability per subset}

\begin{table}[h]
\caption{HIDSAG cross-preprocessing ARI (mean over the 3 recipe
pairs cont-rem / deriv / smooth) per subset. Public artefact
\texttt{hidsag\_cross\_preprocessing\_stability/<subset>.json}.}
\label{tab:b9}
\centering
\small
\begin{tabular}{l c}
\toprule
Subset & mean cross-recipe ARI \\
\midrule
GEOMET   & (reported in artefact) \\
MINERAL1 & (reported in artefact) \\
MINERAL2 & (reported in artefact) \\
GEOCHEM  & (reported in artefact) \\
PORPHYRY & (reported in artefact) \\
\bottomrule
\end{tabular}
\end{table}

\section{Axis F-12: external baseline deltas}

\begin{table}[h]
\caption{Per-scene OA delta between the canonical topic-routed-soft
classifier and three reference papers per scene. Negative values
indicate the baseline outperforms; positive values indicate the
topic-routed-soft outperforms.}
\label{tab:b12}
\centering
\small
\begin{tabular}{l c c c}
\toprule
Scene & vs ref 1 & vs ref 2 & vs ref 3 \\
\midrule
Indian Pines & (reported) & (reported) & (reported) \\
Salinas      & (reported) & (reported) & (reported) \\
Salinas-A    & (reported) & (reported) & (reported) \\
Pavia U      & (reported) & (reported) & (reported) \\
KSC          & (reported) & (reported) & (reported) \\
Botswana     & (reported) & (reported) & (reported) \\
\bottomrule
\end{tabular}
\end{table}

\bibliographystyle{IEEEtran}
\bibliography{../../bibliography/refs}

\end{document}
